\documentclass[11pt,a4paper]{article}
\usepackage[utf8]{inputenc}
\usepackage[margin=1in]{geometry}
\usepackage{amsmath,amssymb,amsthm}
\usepackage{listings}
\usepackage{xcolor}
\usepackage{hyperref}

\newtheorem{theorem}{Theorem}
\newtheorem{lemma}[theorem]{Lemma}

\lstset{
  language=C++,
  basicstyle=\ttfamily\small,
  keywordstyle=\color{blue}\bfseries,
  commentstyle=\color{green!60!black},
  stringstyle=\color{red!70!black},
  numbers=left,
  numberstyle=\tiny\color{gray},
  breaklines=true,
  frame=single,
  tabsize=4,
  showstringspaces=false
}

\title{IOI 2011 -- Hottest (Hot)}
\author{Solution}
\date{}

\begin{document}
\maketitle

\section{Problem Summary}
Given $N$ measurement points in the plane, each at $(x_i, y_i)$ with temperature $t_i$, find a circle of radius $R$ that maximizes the average temperature of enclosed points. Output the center of such a circle.

\section{Solution}

\subsection{Key Observation}
The optimal circle must have at least one measurement point on or near its boundary, or two points defining its boundary. Therefore, the set of candidate centers is:
\begin{enumerate}
    \item Each measurement point itself.
    \item For each pair of points within distance $2R$, the two centers of circles of radius~$R$ passing through both points.
\end{enumerate}

\subsection{Algorithm}
For each candidate center $(c_x, c_y)$, compute the average temperature of all points within distance~$R$. Return the center with the highest average.

For two points $(x_i, y_i)$ and $(x_j, y_j)$ with distance $d \le 2R$, the two candidate centers lie at:
\[
\left(\frac{x_i + x_j}{2} \pm h \cdot \frac{-(y_j - y_i)}{d},\;
      \frac{y_i + y_j}{2} \pm h \cdot \frac{x_j - x_i}{d}\right)
\]
where $h = \sqrt{R^2 - d^2/4}$.

\subsection{Complexity}
\begin{itemize}
    \item \textbf{Time:} $O(N^3)$ --- $O(N^2)$ candidate centers, each evaluated in $O(N)$.
    \item \textbf{Space:} $O(N)$.
\end{itemize}

For larger $N$, an $O(N^2 \log N)$ angular sweep can be used: for each point~$i$, compute the angular intervals of centers that include each other point~$j$ (with $d_{ij} \le 2R$), then sweep to find the maximum weighted count.

\section{Implementation}

\begin{lstlisting}
#include <bits/stdc++.h>
using namespace std;

int main() {
    ios::sync_with_stdio(false);
    cin.tie(nullptr);

    int N;
    double R;
    cin >> N >> R;

    vector<double> x(N), y(N), t(N);
    for (int i = 0; i < N; i++)
        cin >> x[i] >> y[i] >> t[i];

    double bestAvg = -1e18;
    double bestCx = 0, bestCy = 0;

    auto eval = [&](double cx, double cy) {
        double sumT = 0;
        int cnt = 0;
        for (int i = 0; i < N; i++) {
            double dx = x[i] - cx, dy = y[i] - cy;
            if (dx*dx + dy*dy <= R*R + 1e-9) {
                sumT += t[i];
                cnt++;
            }
        }
        if (cnt > 0) {
            double avg = sumT / cnt;
            if (avg > bestAvg) {
                bestAvg = avg;
                bestCx = cx;
                bestCy = cy;
            }
        }
    };

    // Candidate 1: each point as center
    for (int i = 0; i < N; i++)
        eval(x[i], y[i]);

    // Candidate 2: pairs of points on boundary
    for (int i = 0; i < N; i++) {
        for (int j = i + 1; j < N; j++) {
            double dx = x[j] - x[i], dy = y[j] - y[i];
            double d2 = dx*dx + dy*dy;
            if (d2 > 4*R*R + 1e-9) continue;

            double mx = (x[i] + x[j]) / 2;
            double my = (y[i] + y[j]) / 2;
            double h = sqrt(max(0.0, R*R - d2/4));
            double len = sqrt(d2);
            if (len < 1e-12) continue;

            double px = -dy / len, py = dx / len;
            eval(mx + h * px, my + h * py);
            eval(mx - h * px, my - h * py);
        }
    }

    cout << fixed << setprecision(6) << bestCx << " " << bestCy << "\n";
    return 0;
}
\end{lstlisting}

\end{document}
